% Source: https://artofproblemsolving.com/community/c3h3180824_2021_emcc_individual_speed_test__exeter_math_club_competition
% Stable identity: aops-topic-3180824-q17
The divisor count of $n=\prod p_i^{e_i}$ is $\prod(e_i+1)$. It is prime exactly when $n=p^{q-1}$ for a prime $q$. Up to 100 these are the 25 primes themselves, the four prime squares 4, 9, 25, 49, the two fourth powers 16, 81, and the sixth power 64. Higher prime exponents $q-1\ge10$ already give $2^{10}>100$. Also 1 has one divisor, neither prime nor composite. Hence the desired count is $100-25-4-2-1-1=\boxed{67}$.
